% Overview figure (figure*, drawn at \textwidth = 17.46 cm). The Overleaf figure (fig/overview_v1.tex), refined:
% numbered stages with a one-line description each, single-turn inference only (no KV cache), less empty space,
% larger frames and text. Three stages left to right:
% (1) retrieval: video + subtitles -> retriever (with query q) -> relevance r_i over items in time order;
% (2) Token Scheduler (the contribution, the only accented stage): objective, value curves v_frame / v_text,
%     water-filling -> token plan b_i;
% (3) inference: frozen pruner keeps each frame's top-ranked tokens, the frozen VLM answers.
% Bar heights and the plan are illustrative, not measured. Frames: Video-MME video bepwr1-CNRU (illustration only).
% Tableau 10 (Stone, 2016): blue 4E79A7 and orange F28E2B on neutral greys; dark text on light fills only
\definecolor{ink}{HTML}{333333}%
\definecolor{arrowgray}{HTML}{6F6F6F}%
\definecolor{bandA}{HTML}{F4F5F7}%
\definecolor{bandB}{HTML}{EAF1F8}%
\definecolor{edgeB}{HTML}{4E79A7}%
\definecolor{bandC}{HTML}{F4F5F7}%
\definecolor{stepA}{HTML}{333333}%
\definecolor{stepB}{HTML}{4E79A7}%
\definecolor{stepC}{HTML}{333333}%
\definecolor{modelA}{HTML}{DCE6F1}%
\definecolor{modelC}{HTML}{DCE6F1}%
\definecolor{modelD}{HTML}{4E79A7}%
\definecolor{modelT}{HTML}{333333}%
\definecolor{flake}{HTML}{4E79A7}%
\definecolor{chipF}{HTML}{FFFFFF}%
\definecolor{chipD}{HTML}{A0A0A0}%
\definecolor{frm}{HTML}{4E79A7}%
\definecolor{txt}{HTML}{F28E2B}%
\definecolor{accA}{HTML}{4E79A7}%
% every color role: ink, arrowgray, bandA/B/C, edgeB, stepA/B/C, modelA/C/D/T, flake, chipF/D, frm, txt
\newcommand{\snowflake}[2]{%
  \begin{scope}[shift={#1}, line width=0.55pt, draw=flake, line cap=round]
    \foreach \a in {0,60,...,300}{\draw (0,0) -- (\a:#2);
      \draw (\a:0.55*#2) -- ++(\a+45:0.32*#2); \draw (\a:0.55*#2) -- ++(\a-45:0.32*#2);}
  \end{scope}}
% \items{width}{height scale}{v1/t1,...,v20/t20}: 20 items on one time axis; t = 1 marks a transcript event (light bar)
\newcommand{\items}[3]{%
  \foreach \v/\t [count=\i] in {#3} {
    \ifnum\t=1 \fill[txt] ({(\i-0.5)*#1/20-0.3*#1/20},0) rectangle ++({0.6*#1/20},{\v*#2});
    \else      \fill[frm]    ({(\i-0.5)*#1/20-0.3*#1/20},0) rectangle ++({0.6*#1/20},{\v*#2}); \fi}}
% \prunedframe{center}{image}{kept tokens out of 32}: a frame (1.0 cm) as an 8x4 token grid; dropped tokens greyed out.
% The keep order is one fixed ranking (centre first), so a frame at b tokens holds its top b tokens.
\newcommand{\prunedframe}[3]{%
  \node[inner sep=0pt] at (#1) {\includegraphics[width=1.0cm]{fig/frames/#2}};
  \begin{scope}[shift={($(#1)+(-0.5,-0.28125)$)}]
    \foreach \c [count=\k from 0] in {10,11,18,19,9,12,17,20,2,3,26,27,1,4,25,28,8,13,16,21,0,5,24,29,6,14,22,30,7,15,23,31} {
      \ifnum\k<#3\relax\else
        \pgfmathtruncatemacro{\cx}{mod(\c,8)}\pgfmathtruncatemacro{\cy}{int(\c/8)}
        \fill[white, opacity=0.85] (\cx*0.125,\cy*0.140625) rectangle ++(0.125,0.140625);
      \fi}
    \draw[ink!30, line width=0.15pt, xstep=0.125, ystep=0.140625] (0,0) grid (1.0,0.5625);
    \draw[ink!45, line width=0.3pt] (0,0) rectangle (1.0,0.5625);
  \end{scope}}
\newcommand{\txtchip}[1]{% transcript event: a short text chip
  \fill[txt!35, rounded corners=1pt] ($(#1)+(-0.42,-0.2)$) rectangle ++(0.84,0.4);
  \draw[txt!85!black, line width=0.5pt] ($(#1)+(-0.3,0.07)$) -- ++(0.6,0);
  \draw[txt!85!black, line width=0.5pt] ($(#1)+(-0.3,-0.07)$) -- ++(0.4,0);}
% \stage{x}{y}{number}{title}{one-line description}: numbered stage heading
\newcommand{\stage}[5]{%
  \node[circle, fill=#5, text=white, font=\sffamily\scriptsize\bfseries, inner sep=0pt, minimum size=0.4cm] at (#1+0.2,#2) {#3};
  \node[font=\sffamily\small\bfseries, text=#5, anchor=west, inner sep=0pt] at (#1+0.5,#2) {#4};}

\begin{tikzpicture}[font=\sffamily\footnotesize, text=ink, line width=0.6pt,
  card/.style={fill=white, draw=ink!18, line width=0.5pt, rounded corners=4pt, align=center, inner sep=3.5pt},
  arr/.style={-{Stealth[length=5pt, width=4pt]}, draw=arrowgray, line width=0.8pt},
  ax/.style={-{Stealth[length=3.5pt, width=3pt]}, draw=ink!55, line width=0.5pt},
  axl/.style={font=\sffamily\scriptsize, text=ink!70},
  desc/.style={font=\sffamily\scriptsize, text=ink!65, anchor=north west, inner sep=0pt, align=left},
  elbl/.style={font=\sffamily\scriptsize, text=ink!75, inner sep=1.5pt},
  model/.style={fill=modelC, draw=modelD, text=modelT, line width=0.5pt, rounded corners=3pt, align=center,
    minimum height=0.62cm, inner xsep=5pt, font=\sffamily\footnotesize},
  qa/.style={fill=chipF, draw=chipD, line width=0.5pt, rounded corners=3pt, align=left,
    inner sep=3pt, font=\sffamily\scriptsize}]

% ---------------- layout ----------------
% Top row: (1) Retrieval -> (2) Allocation; bottom row: (3) Inference, left to right. Stages connect box to box
% (r_i, b_i). Q sits inside Retrieval and A inside Inference; nothing is outside the stage boxes.
\fill[rounded corners=6pt, bandA] (0,2.0) rectangle (7.5,5.3);
\fill[rounded corners=6pt, bandB] (8.0,2.0) rectangle (17.4,5.3);
\draw[rounded corners=6pt, edgeB, line width=0.6pt] (8.0,2.0) rectangle (17.4,5.3);
\fill[rounded corners=6pt, bandC] (0,0) rectangle (17.4,1.6);
\stage{0.12}{5.02}{1}{Retrieval}{stepA}
\stage{8.12}{5.02}{2}{Allocation}{stepB}
\stage{0.12}{1.3}{3}{Inference}{stepC}
\draw[arr] (7.5,3.65) -- node[elbl, above] {$r_i$} (8.0,3.65);
\draw[arr] (12.7,2.0) -- node[elbl, right] {$b_i$} (12.7,1.6);

% ---------------- (1) retrieval ----------------
\fill[ink!88, rounded corners=1pt] (0.2,3.85) rectangle (7.3,4.75);
\foreach \x in {0.28,0.48,...,7.2} {\fill[white] (\x,4.65) rectangle ++(0.1,0.06); \fill[white] (\x,3.89) rectangle ++(0.1,0.06);}
\foreach \k in {1,...,5} {\node[inner sep=0pt] at ({0.91+(\k-1)*1.42},4.3) {\includegraphics[width=1.2cm]{fig/frames/frame\k.jpg}};}
\fill[txt!35, rounded corners=1pt] (0.2,3.57) rectangle (7.3,3.77);
\foreach \x/\l in {0.5/1.0, 1.9/0.8, 3.1/1.2, 4.7/0.9, 6.0/0.9} {\draw[txt!85!black, line width=0.5pt] (\x,3.67) -- ++(\l,0);}
\node[qa] (q) at (1.0,2.75) {\textbf{Q:} Color of\\the last car?};
\node[model, fill=modelA] (ret) at (2.95,2.75) {Retriever\hspace{0.12cm}};
\snowflake{($(ret.north east)+(-0.14,-0.14)$)}{0.075}
\draw[arr] (q.east) -- (ret.west);
\draw[arr] (2.95,3.55) -- (ret.north);
\begin{scope}[shift={(4.25,2.42)}]
  \draw[ax] (0,0) -- (3.05,0); \draw[ax] (0,0) -- (0,1.0);
  \node[axl, anchor=west] at (0.06,0.93) {$r_i$};
  \node[axl, anchor=north] at (1.5,-0.02) {frame \#};
  \items{2.95}{1.4}{.04/0,.06/0,.30/0,.52/0,.20/1,.40/0,.10/0,.05/0,.04/0,.36/1,
                   .06/0,.22/0,.30/0,.42/1,.34/0,.26/0,.08/0,.18/1,.05/0,.03/0}
\end{scope}
\draw[arr] (ret.east) -- (4.15,2.75);

% ---------------- (2) allocation ----------------
\node[font=\sffamily\scriptsize] at (12.7,4.5) {$\max\limits_{\mathbf b}\ \sum_i r_iv_i(b_i)-\eta\sum_{i\neq j}s_{ij}\,v_i(b_i)\,v_j(b_j)\qquad\text{s.t.}\ \ \sum_i b_i\le B$};
\begin{scope}[shift={(8.75,2.42)}]
  \draw[ax] (0,0) -- (3.2,0); \draw[ax] (0,0) -- (0,1.6);
  \node[axl, anchor=west] at (0.06,1.5) {$v(b)$};
  \node[axl, anchor=north] at (1.6,-0.02) {$b$};
  \draw[txt, line width=1.1pt, domain=0.02:8, samples=60] plot ({0.37*\x}, {1.45*((2^\x-1)^0.841)/(((2^\x-1)^0.841)+3.44)});
  \draw[frm, line width=1.1pt, domain=0.02:8, samples=60] plot ({0.37*\x}, {1.35*((2^\x-1)^0.841)/(((2^\x-1)^0.841)+7.73)});
  \node[font=\sffamily\scriptsize, text=txt!80!black, anchor=south west, inner sep=1pt] at (3.0,1.35) {$v_{\text{text}}$};
  \node[font=\sffamily\scriptsize, text=frm, anchor=north west, inner sep=1pt] at (3.0,1.2) {$v_{\text{frame}}$};
\end{scope}
\draw[arr] (12.75,3.2) -- (13.45,3.2);
\begin{scope}[shift={(13.85,2.42)}]
  \draw[ax] (0,0) -- (3.3,0); \draw[ax] (0,0) -- (0,1.6);
  \node[axl, anchor=west] at (0.06,1.5) {$b_i$};
  \node[axl, anchor=north] at (1.6,-0.02) {frame \#};
  \items{3.2}{1.9}{0/0,0/0,.438/0,.688/0,.25/1,.531/0,0/0,0/0,0/0,.5/1,
                   0/0,.25/0,.344/0,.5/1,.406/0,.312/0,0/0,.25/1,0/0,0/0}
\end{scope}

% ---------------- (3) inference ----------------
% the frozen pruner is a container that holds its output: the frames at their token counts and the text events
\node[model, minimum width=9.4cm, minimum height=1.0cm] (pr) at (7.3,0.7) {};
\node[font=\sffamily\footnotesize, text=modelT, align=center] at (3.4,0.7) {Token\\pruner};
\snowflake{($(pr.north east)+(-0.14,-0.14)$)}{0.075}
\prunedframe{4.95,0.7}{frame1.jpg}{6}
\txtchip{6.15,0.7}
\prunedframe{7.35,0.7}{frame3.jpg}{22}
\prunedframe{8.55,0.7}{frame4.jpg}{12}
\txtchip{9.75,0.7}
\prunedframe{10.95,0.7}{frame5.jpg}{32}
\node[model, minimum width=2.4cm, minimum height=1.0cm] (vlm) at (14.0,0.7) {VLM};
\snowflake{($(vlm.north east)+(-0.14,-0.14)$)}{0.075}
\node[qa] (ans) at (16.4,0.7) {\textbf{A:} Blue};
\draw[arr] (pr.east) -- (vlm.west);
\draw[arr] (vlm.east) -- (ans.west);
\end{tikzpicture}
